\chapter{Kết luận}
\label{Chapter6}

\section{Các kết quả đã đạt được của hệ thống ứng dụng}

Luận văn này đã nghiên cứu, phân tích và hiện thực hóa thành công hệ thống CareerNova --- một nền tảng hỗ trợ sinh viên ngành Công nghệ thông tin đánh giá khoảng cách năng lực và khám phá thị trường lao động thông qua trực quan hóa dữ liệu. Kết thúc quá trình xây dựng và thử nghiệm, đề tài đã đạt được các kết quả chính sau đây:

\begin{enumerate}
    \item \textbf{Xây dựng hoàn chỉnh khung dữ liệu đa nguồn}: Phát triển thành công hệ thống thu thập dữ liệu tin tuyển dụng tự động từ bốn nền tảng lớn tại Việt Nam (ITViec, VietnamWorks, CareerViet, LinkedIn) theo cơ chế dự phòng hai tầng (cào tĩnh ưu tiên, cào động dự phòng), cài đặt các module tiền xử lý dữ liệu và thuật toán khử trùng lặp trước khi lưu trữ tập trung. Số lượng dữ liệu tích lũy và tin đang mở được ghi nhận theo snapshot cơ sở dữ liệu tại thời điểm kiểm thử trong Mục 5.3.1.

    \item \textbf{Ứng dụng thành công công nghệ AI tiên tiến}: Áp dụng mô hình ngôn ngữ lớn Gemini để trích xuất thông tin kỹ năng và vị trí công việc từ văn bản phi cấu trúc thông qua kỹ thuật Prompt Engineering có cấu trúc (Structured Outputs kết hợp JSON Schema) và cơ chế hậu kiểm chống ảo giác dựa trên bằng chứng. Kết quả thực nghiệm trên tập dữ liệu đối soát cho thấy module trích xuất đạt Precision 98.2\%, Recall 97.8\% và F1-Score 98.0\%, vượt xa ngưỡng kỳ vọng 82.5\% đặt ra ban đầu.

    \item \textbf{Đề xuất giải thuật đối soát và phân tích khoảng cách kỹ năng thông minh}: Kết hợp so khớp ngữ nghĩa bằng vector nhúng SBERT cục bộ (\texttt{all-MiniLM-L6-v2}) với trọng số TF-IDF được tính động từ dữ liệu thị trường thực tế và bộ sáu quy tắc lọc ngữ nghĩa theo metadata phân loại kỹ năng, giúp định lượng hóa chính xác mức độ tương thích giữa hồ sơ CV của ứng viên và yêu cầu tuyển dụng, đồng thời phân loại kỹ năng thành ba nhóm hành động được (đã có, cần bổ sung, thiếu hoàn toàn).

    \item \textbf{Hiện thực hóa và triển khai ứng dụng Web hoàn chỉnh}: Triển khai hệ thống thực tế tại địa chỉ \texttt{career-nova.online} trên hạ tầng đám mây, cung cấp đầy đủ các phân hệ: xác thực người dùng (email và Google SSO), Dashboard trực quan hóa xu hướng thị trường lao động, quản lý CV, tìm kiếm và lưu việc làm, phân tích đối soát CV với biểu đồ Radar, báo cáo khoảng cách kỹ năng và gợi ý lộ trình học tập cá nhân hóa theo kỹ năng thiếu hụt. Kết quả kiểm thử chức năng và số liệu hiệu năng được tổng hợp tại Chương 5.
\end{enumerate}

Về mặt học thuật, đề tài đóng góp vào kho tri thức về ứng dụng mô hình ngôn ngữ lớn trong lĩnh vực phân tích thị trường lao động, đặc biệt là phương pháp kết hợp Structured Outputs với cơ chế kiểm chứng chống ảo giác và kiến trúc lai LLM + SBERT nhằm cân bằng giữa chất lượng ngữ nghĩa và hiệu quả kinh tế trong bài toán so khớp hồ sơ nhân sự.

%==========================================================================

\section{Hạn chế còn tồn tại của đề tài}

Mặc dù hệ thống đã vận hành ổn định và đạt các chỉ số đánh giá tốt trong giai đoạn thử nghiệm, đề tài vẫn còn tồn tại một số hạn chế cần được nhìn nhận khách quan:

\begin{itemize}
    \item \textbf{Phạm vi dữ liệu}: Thử nghiệm thực tế mới chỉ tập trung chủ yếu vào nhóm ngành Công nghệ Thông tin (IT) và Khoa học Dữ liệu tại thị trường Việt Nam, chưa mở rộng sang các lĩnh vực ngành nghề phi kỹ thuật khác. Quy mô dữ liệu tích lũy đủ phục vụ phân tích xu hướng ngắn hạn trong phạm vi triển khai, nhưng chưa đủ dày để phân tích chu kỳ tuyển dụng theo mùa qua nhiều năm.

    \item \textbf{Hiệu năng thời gian thực}: Thời gian phản hồi của tính năng đối soát và phân tích khoảng cách kỹ năng vẫn còn phụ thuộc vào độ trễ của API mô hình ngôn ngữ lớn bên ngoài (trong trường hợp cache miss) và thời điểm hệ thống đồng bộ dữ liệu tuyển dụng. Hai endpoint thống kê kỹ năng của Dashboard (\texttt{in-demand}, \texttt{rising}) chậm hơn đáng kể các endpoint còn lại do truy vấn tổng hợp trên bảng liên kết lớn, cần tối ưu bằng chỉ mục bổ sung hoặc materialized view.

    \item \textbf{Sự phụ thuộc vào chất lượng dữ liệu thô}: Kết quả trích xuất kỹ năng có thể bị ảnh hưởng nếu tệp CV của ứng viên được trình bày quá cẩu thả, ảnh quét có độ phân giải kém hoặc sử dụng các định dạng thiết kế cột phức tạp mà tầng OCR không đọc đúng thứ tự dòng.

    \item \textbf{Từ điển kỹ năng cần bảo trì định kỳ}: Các kỹ năng mới xuất hiện trên thị trường chưa có trong từ điển chuẩn hóa được đẩy vào bảng \texttt{unmatched\_skills} và hiện vẫn cần con người phê duyệt thủ công trước khi bổ sung vào từ điển, chưa có quy trình tự động hóa hoàn chỉnh.
\end{itemize}

%==========================================================================

\section{Bài học kinh nghiệm và Kỹ năng tích lũy}

Trải qua toàn bộ vòng đời phát triển của dự án từ khâu khảo sát nghiệp vụ, phân tích thiết kế, đến triển khai mã nguồn và kiểm thử vận hành, nhóm thực hiện đề tài đã tích lũy được những bài học kinh nghiệm và kỹ năng thực tiễn quý giá, làm hành trang vững chắc cho con đường kỹ sư phần mềm chuyên nghiệp:

\begin{enumerate}
    \item \textbf{Tư duy thiết kế hệ thống phân tán (System Architecture)}: Nhóm học được cách phá vỡ một ứng dụng nguyên khối khổng lồ thành các dịch vụ độc lập (Frontend Web, Backend Core API, ETL Data Pipeline). Việc áp dụng container hóa bằng Docker và Docker Compose ngay từ đầu giúp nhóm thấu hiểu nguyên lý triển khai nhất quán giữa môi trường phát triển (dev) và môi trường sản phẩm (production), giảm thiểu triệt để lỗi "chạy được trên máy tôi nhưng lỗi trên server".

    \item \textbf{Làm chủ chu trình xử lý dữ liệu (Data Engineering)}: Thông qua việc xây dựng ETL Pipeline, nhóm nắm vững kỹ thuật thu thập dữ liệu tự động, cách vượt rào cản chống bot bằng cơ chế dự phòng hai tầng (Request tĩnh kết hợp Selenium headless), cũng như các kỹ thuật làm sạch, chuẩn hóa và tối ưu truy vấn dữ liệu lớn (sử dụng JSONB thay vì JOIN chéo) theo mô hình phân tích OLAP.

    \item \textbf{Tích hợp AI vào sản phẩm thực tế (Applied AI \& LLMOps)}: Nhóm nhận ra rằng việc ứng dụng AI trong thực tế phức tạp hơn nhiều so với lý thuyết. Bài học lớn nhất là cách "kìm cương" mô hình ngôn ngữ lớn (LLM) thông qua kỹ thuật \textit{Structured Outputs} và cơ chế kiểm chứng chéo (Anti-Hallucination) để ép AI sinh dữ liệu có cấu trúc ổn định thay vì trả lời tự do. Nhóm cũng thấu hiểu sự cân bằng giữa hiệu năng và chi phí khi kết hợp LLM (nhạy bén nhưng chậm, tốn phí) với Vector Search/SBERT cục bộ (nhanh, miễn phí, tính toán toán học).

    \item \textbf{Tư duy kiểm thử và tư duy định lượng (QA \& Quantitative Mindset)}: Trong giai đoạn đánh giá hệ thống, nhóm rèn luyện được thói quen làm việc với các con số (Data-driven). Việc sử dụng Apache JMeter để ép tải hệ thống và đo lường độ trễ API tính bằng mili-giây giúp nhóm gạt bỏ những đánh giá cảm tính, nhìn nhận thẳng thắn vào các điểm thắt nút cổ chai (bottleneck) của kiến trúc, từ đó hình thành tư duy liên tục tối ưu hóa hiệu năng phần mềm.
\end{enumerate}

%==========================================================================

\section{Hướng nghiên cứu và phát triển tương lai}

Dựa trên nền tảng kỹ thuật và dữ liệu đã xây dựng, cùng với việc nhìn nhận các hạn chế hiện tại, các hướng nghiên cứu và phát triển có giá trị cao trong tương lai bao gồm:

\begin{enumerate}
    \item \textbf{Mở rộng đa ngành nghề}: Cập nhật cây phân loại (taxonomy) và huấn luyện thêm các mô hình nhúng để hỗ trợ phân tích dữ liệu tuyển dụng ở các lĩnh vực khác như Tài chính, Marketing, Y tế và Giáo dục, đồng thời mở rộng danh sách nguồn thu thập để tăng độ phủ thị trường.

    \item \textbf{Tự động hóa vòng lặp cải thiện từ điển kỹ năng}: Xây dựng quy trình bán tự động khai thác bảng \texttt{unmatched\_skills} --- kết hợp tần suất xuất hiện, điểm tương đồng gần nhất và xác nhận của LLM --- để đề xuất bổ sung kỹ năng mới vào từ điển chuẩn hóa theo chu kỳ, giúp hệ thống tự thích nghi với các công nghệ mới nổi trên thị trường mà không cần can thiệp thủ công.

    \item \textbf{Nâng cao chất lượng gợi ý lộ trình học tập}: Mở rộng kho khóa học liên kết với Gap Report bằng cách tự động đồng bộ metadata từ các nền tảng học trực tuyến (Coursera, Udemy, edX), bổ sung cơ chế xếp hạng khóa học theo phản hồi người dùng và theo dõi tiến độ học tập để đo lường mức độ thu hẹp khoảng cách kỹ năng theo thời gian.

    \item \textbf{Tối ưu hóa hiệu năng và triển khai mô hình nội bộ (On-premise LLMs)}: Nghiên cứu tinh chỉnh (Fine-tuning) và triển khai các mô hình ngôn ngữ lớn mã nguồn mở có kích thước nhỏ (như LLaMA-3-8B hoặc Qwen-2.5) trực tiếp trên máy chủ hệ thống để tự chủ công nghệ, tăng tính bảo mật dữ liệu và giảm tối đa thời gian trễ của khâu trích xuất kỹ năng khi cache miss.
\end{enumerate}

Nhìn chung, hệ thống CareerNova đã xây dựng được một nền tảng dữ liệu và kỹ thuật vững chắc, có khả năng mở rộng sang các hướng ứng dụng thực tiễn rộng lớn hơn trong lĩnh vực hỗ trợ quyết định nghề nghiệp và phân tích thị trường lao động bằng trí tuệ nhân tạo. Đề tài khẳng định tính khả thi của việc kết hợp các mô hình ngôn ngữ lớn thế hệ mới với kỹ thuật biểu diễn vector ngữ nghĩa truyền thống để giải quyết bài toán thực tiễn có ý nghĩa giáo dục và xã hội sâu sắc: thu hẹp khoảng cách thông tin giữa sinh viên và thị trường lao động, trao quyền cho thế hệ trẻ chủ động định hướng con đường nghề nghiệp của bản thân dựa trên dữ liệu.
